\documentclass[11pt]{article}
\usepackage[margin=1.1in]{geometry}
\usepackage{amsmath,amssymb}
\usepackage{booktabs}
\usepackage{array}
\usepackage{graphicx}
\usepackage[numbers,sort&compress]{natbib}
\usepackage[colorlinks=true,linkcolor=blue,citecolor=blue,urlcolor=blue]{hyperref}
\usepackage{microtype}

\newcolumntype{L}[1]{>{\raggedright\arraybackslash}p{#1}}
\setlength{\belowcaptionskip}{7pt}
\renewcommand{\arraystretch}{1.15}

\title{If two genes look interchangeable, are they?\\
A calibration of expression redundancy against\\
interventional equivalence in 43 million gene pairs}
\author{Akanda Wahid-Ul-Ashraf}
\date{August 2026}

\begin{document}
\maketitle

\begin{abstract}
A routine step in computational biology treats two genes that track each
other closely in expression data as one hypothesis: pick either, and the
other would have behaved the same. We could find no published estimate of
how often that holds, so we measured it. Using DepMap 24Q4---1,103 cancer
cell lines with paired expression profiles and genome-wide CRISPR knockout
screens---we scored 43.9 million gene pairs on two axes: how redundant the
two genes look observationally, and how similar the consequences of knocking
each one out actually are. Redundancy is real evidence and it is far from
identification. Relative to a base rate of 0.42\% (95\% CI 0.37--0.47), the
odds of interventional equivalence rise about 39-fold (95\% CI 23--65) in
the most redundant well-populated bins, but the ceiling there is 17.0\%
(95\% CI 9.7--27.3); among the 13 most redundant pairs in the dataset, none
were equivalent. Three secondary results follow. Cell lineage manufactures
apparent redundancy: correcting for it removes 56\% of pairs at $r^2 \ge
0.6$. Panel redundancy---predictability from ten partners rather than
one---carries essentially the same information as best-pair redundancy
(AUC 0.614 versus 0.607), so the simpler statistic loses nothing. And genes
inside a real gene family show a \emph{lower} equivalence ceiling (13.9\%)
than genes outside one (26.3\%), consistent with the buffering that makes
single-knockout screens structurally blind to compensating pairs. We report
two pre-registered failures at full volume: a null control that could not
fail, and a positive control whose failure correctly indicted our first
similarity measure rather than the biology. Code, pre-registration and the
complete experiment ledger are public.
\end{abstract}

\section{The assumption}

Given expression profiles across many samples, two genes sometimes track
each other so closely that they appear interchangeable. A common next step
is to treat them as a single hypothesis---to perturb one, target one, or
follow one up, on the understanding that the other would have behaved
similarly. The step is rarely stated explicitly, which is part of why it is
rarely checked.

Stated explicitly it is a conditional probability. Let $r_{\mathrm{obs}}$
measure how redundant two genes look observationally, and let a pair be
\emph{interventionally equivalent} when knocking out either produces the
same phenotype. The quantity the everyday decision relies on is
\[
  P\!\left(\text{equivalent under intervention} \mid r_{\mathrm{obs}} = r\right),
\]
as a function of $r$. This paper measures that function.

\section{What has and has not been measured}

Adjacent literature is substantial and should not be mistaken for an answer
to this question.

Co-essentiality analyses ask the \emph{reverse} direction: given that two
genes have similar knockout profiles, what else do they share? That is
informative about pathway membership but says nothing about how often
observational similarity licenses an interventional claim, which is the
direction in which decisions are actually made.

Closer still is the work of \citet{dekegel2019} on paralog buffering, and
its extension to genome-scale synthetic-lethality prediction
\citep{dekegel2021,dekegel2026}. Those studies ask whether losing one
member of a paralog pair \emph{increases dependency} on the other. That is
a question about compensation---about pairs whose members cover for one
another---and it is measured on paralogs specifically. Interventional
equivalence is a different relation: it asks whether the two knockouts
produce the \emph{same} outcome, over the whole gene universe rather than
paralogs alone. The two are not merely different, they point opposite ways:
a perfectly buffering pair is maximally non-equivalent under single
knockout, precisely because each gene covers for the other. Section
\ref{sec:families} shows this predicted opposition in our own data, which we
take as convergent evidence for their finding rather than a new one.

We could not find the forward calibration published. We state that as the
result of a search rather than as a claim about the literature, and we would
welcome a pointer to prior work.

\section{Data and measures}

\paragraph{Data.} DepMap 24Q4: expression (\texttt{OmicsExpression\-Protein\-Coding\-GenesTPMLogp1})
and CRISPR gene effect (Chronos, \citealp{dempster2021}) for the
1,103 cell lines present in both, restricted to the 17,716 genes present in
both matrices. All files are public and downloaded by the analysis script;
none are redistributed.

\paragraph{Gene universe.} Genes are retained if mean expression clears a
floor and expression variance clears the 10th percentile of expressed genes,
leaving 9,368 genes and 43,875,028 unordered pairs. Filters were fixed in
the pre-registration before any curve was computed.

\paragraph{Observational axis.} $r_{\mathrm{obs}}$ is the squared Pearson
correlation of expression across cell lines, after regressing out one-hot
cell lineage (\S\ref{sec:lineage} reports the uncorrected arm).

\paragraph{Interventional axis.} Our first declared measure was the Pearson
correlation of Chronos profiles, $e_{\mathrm{sel}}$. It failed its positive
control (\S\ref{sec:controls}), and the pre-registration contained a rule
for that case. The reported measure is proximity on \emph{raw, uncentred}
Chronos vectors,
\[
  e_{\mathrm{prox}}(a,b) \;=\;
  \frac{2\langle k_a, k_b\rangle}{\langle k_a,k_a\rangle + \langle k_b,k_b\rangle},
\]
which equals 1 for identical profiles and stays near 0 for unrelated ones.
Centring, which Pearson performs, removes exactly the shared mean in which
uniform co-essentiality lives; proximity keeps it. A pair counts as
interventionally equivalent when $e_{\mathrm{prox}} > 0.8$. Both axes are
reported throughout, and the qualitative shape of the curve is unchanged by
moving the threshold; the numerical ceiling is not, and depends on it.

\paragraph{Uncertainty.} Gene pairs share genes, so pairs are not
independent. All intervals are gene-level bootstrap, $B=100$, resampling
genes and recomputing the entire curve.

\section{Controls, including the two that failed}
\label{sec:controls}

\paragraph{A null that could not fail.} The pre-registered null permuted the
rows (cell lines) of one matrix. Its output was bin-for-bin identical to the
real analysis---which is how the error surfaced. Within-matrix correlations
are invariant under a shared row permutation, so the control was
algebraically incapable of differing from the thing it was meant to test. It
was replaced with a gene-label permutation, which can fail. Under the
replacement the calibration curve flattens: base rate 0.422\% and a maximum
bin equivalence of 0.9\%, against 17\% in the real data. \emph{A control that
cannot fail is not a control.}

\paragraph{A positive control that failed, correctly.} 253 same-complex gene
pairs drawn from CORUM are the most functionally equivalent pairs biology
offers. Under $e_{\mathrm{sel}}$ they scored a median of 0.221, with
\emph{none} above the equivalence threshold; PSMA1--PSMB5, two subunits of
the same proteasome, scored 0.04. The pre-registration directed that a
failing positive control indicts the measure before the biology. It did:
these genes are near-uniformly essential, so their Chronos profiles are flat
and highly correlated in the mean that Pearson discards. Under
$e_{\mathrm{prox}}$ the same pairs score a median of 0.925 with 88.1\% above
threshold, and PSMA1--PSMB5 scores 0.90.

\paragraph{A matched negative control.} 1,999 randomly drawn gene pairs
score a median $e_{\mathrm{prox}}$ of 0.026, with 0.55\% above
threshold---indistinguishable from the 0.42\% base rate, as required.

\section{Result: the calibration}

\begin{figure}[t]
\centering
\includegraphics[width=0.92\linewidth]{../ExpOutput/depmap_calibration/fig_calibration_cover.png}
\caption{Probability that two genes are interventionally equivalent, as a
function of how redundant they look in expression. Band is the 95\%
gene-level bootstrap interval; the dashed line is the base rate over all
43.9M pairs. The isolated point pools $r^2 \ge 0.8$, where 0 of 13 pairs are
equivalent; its whisker is the rule-of-three upper bound, not a bootstrap
interval, because a bootstrap of an all-zero count is degenerate.}
\label{fig:curve}
\end{figure}

Figure \ref{fig:curve} and Table \ref{tab:headline} give the result.
Redundancy is genuine evidence: the probability of interventional
equivalence rises monotonically across the populated range, reaching
$39\times$ the base rate. It is also nowhere near identification: the
ceiling is about one in six, and the trend does not continue into the
extreme. Among the 13 most redundant pairs in the entire dataset, none are
equivalent---a count small enough that it excludes only rates above roughly
one in four, but one that gives no support to the idea that the curve keeps
climbing.

\begin{table}[t]
\centering
\begin{tabular}{L{6.4cm} r l}
\toprule
Quantity & Estimate & 95\% CI \\
\midrule
Base rate, all 43.9M pairs & 0.42\% & 0.37--0.47 \\
Ceiling, $r^2 \in [0.60, 0.70)$ (911 pairs) & 17.0\% & 9.7--27.3 \\
Lift over base rate & $39.3\times$ & 23.0--64.8 \\
Extreme redundancy, $r^2 \ge 0.8$ & 0 of 13 & $\le 23\%$ \\
\bottomrule
\end{tabular}
\caption{Headline calibration. Intervals are gene-level bootstrap,
$B = 100$; the last row is a rule-of-three bound.}
\label{tab:headline}
\end{table}

For a practitioner the operational statement is: expression redundancy
should raise suspicion that two genes are one hypothesis and should never
settle it. At the redundancy levels where such judgements are typically
made, four times in five or worse, the gene not chosen behaves differently.

\section{Tissue identity manufactures redundancy}
\label{sec:lineage}

Correcting for cell lineage removes 56\% of pairs at $r^2 \ge 0.6$
(2,317 $\rightarrow$ 1,023), and the fraction removed grows with redundancy:
37\% at $r^2 \ge 0.4$, 42\% at $\ge 0.5$, about 70\% at $\ge 0.7$. Two genes
that appear to track each other frequently track the tissue instead. Any
redundancy screen run on a multi-lineage panel without this correction is
measuring cell-of-origin in substantial part, and increasingly so at exactly
the redundancy levels that look most convincing.

\section{Panel redundancy adds nothing over best-pair}

The calibration asks about one partner. A natural objection is that real
redundancy is many-to-one: a gene may be predictable from a \emph{panel} of
others without any single partner standing out. We tested it. For each gene
we computed best-pair $r_{\mathrm{obs}}$, panel $R^2$ from a ridge fit on its
ten most correlated partners, and whether any partner is interventionally
equivalent. The base rate of a gene having some equivalent partner is 0.542.
Panel redundancy predicts that outcome with AUC 0.614; best-pair redundancy
with AUC 0.607. At a matched cut of 0.6, best-pair gives $P = 0.735$
($n = 623$) and panel gives $P = 0.677$ ($n = 2{,}687$). The
pre-registered question is answered in the negative: the more elaborate
statistic carries essentially the same information, and the simpler one is
no worse.

\section{Gene families show the opposite sign, as buffering predicts}
\label{sec:families}

Our pre-registered paralog flag was a proxy---a shared alphabetic symbol
root. Replacing it with HGNC gene groups shows the proxy agrees only 70.3\%
of the time (3,644 false positives, 1,626 missed), which is worth knowing
for any study that uses such a proxy. The published direction is unchanged:
excluding real families, the ceiling is 26.3\%, against 23.1\% under the
proxy.

The new observation is the contrast between family and non-family genes.
Genes \emph{inside} a real gene family have a lower equivalence ceiling
(13.9\%, 388 pairs) than genes outside one (26.3\%, 167 pairs). This is the
sign that buffering predicts: paralogs that cover for one another are, under
\emph{single} knockout, measured as non-equivalent precisely because the
compensation hides the phenotype. It is consistent with
\citet{dekegel2019,dekegel2021} and we report it as convergence with their
result rather than as an independent finding. It also marks the blind spot
of this study honestly: for buffering pairs, our measure returns
non-equivalence for a reason that has nothing to do with the genes being
functionally different.

\section{The subgroup analysis that failed, as declared}

The pre-registration included a triple-negative breast cancer arm and
declared in advance that it would be underpowered. It was, and we report it
as such rather than rescuing it.

Two things went wrong before any result was produced, both instructive. The
first receptor-low definition---lower tertile of ESR1, PGR and ERBB2 among
breast lines---selected zero lines, because PGR's lower tertile \emph{is}
zero and a strict inequality against a threshold on the distribution's floor
excludes exactly the receptor-negative lines it is meant to select. With an
inclusive bound it selected 7. We then gated the definition itself against
documented cell-line receptor status, using no dependency data: it recovered
4 of 19 known triple-negative lines. The definition failed its positive
control, so the definition was replaced, not the biology reinterpreted. The
replacement---ESR1 and PGR at or below the conventional not-expressed line,
ERBB2 at or below the breast-panel median---recovers 10 of 19 and admits
none of the 8 documented receptor-positive lines.

On the resulting 18 lines the calibration degenerates. 11.3\% of pairs land
above $r^2 = 0.6$, against 0.3\% in the 50-line breast panel: with 18
samples, $r^2$ is a noisy statistic and the observational axis inflates
wholesale. Lift collapses from $11\times$ (breast panel) to $3.3\times$. We
draw no subgroup conclusion. The generalisable point is that a subgroup
calibration is not a small version of a pan-cohort one; below a few dozen
samples the redundancy axis stops carrying the information the curve
requires.

\section{Limits}

Cell lines are not tumours, and knockout growth effect in culture is not
clinical response. Nothing here nominates a target or supports a clinical
claim. Single-gene knockouts are structurally blind to buffering, and
\S\ref{sec:families} quantifies rather than removes that blindness. The
numerical ceiling depends on the equivalence threshold; the shape of the
curve does not. The extreme-redundancy result rests on 13 pairs and should
be read as an absence of support for continued climbing, not as a
demonstration of zero. Chronos scores carry their own preprocessing
assumptions, which we inherit.

\section{Conclusion}

Expression redundancy is evidence about interventional equivalence, worth
roughly a 39-fold change in odds, and it is not identification: at the
redundancy levels where the judgement is usually made, the probability that
two interchangeable-looking genes are actually interchangeable is about one
in six at best. Lineage correction is not optional, panel redundancy buys
nothing over the simplest statistic, and gene-family membership moves the
answer in the direction compensation predicts. The number is small enough to
change how a shortlist should be read, and it was cheap to obtain: the whole
calibration runs in about a minute on a consumer GPU against a public
dataset, and can be recomputed for any subtype, panel or release by anyone
who wants a context-specific version.

\section*{Data, code and pre-registration}

All data are public DepMap releases, downloaded by the analysis scripts and
not redistributed. The pre-registration fixing filters, measures, thresholds
and controls was written before any curve was computed. The complete
experiment ledger, including every voided result and both control failures
described above, is published alongside the code. Every number in this paper
is re-derived from its output file by an audit script.

\begin{thebibliography}{9}

\bibitem{dekegel2019} B.~De~Kegel, C.~J.~Ryan.
Paralog buffering contributes to the variable essentiality of genes in
cancer cell lines. \emph{PLOS Genetics} 15(10):e1008466, 2019.

\bibitem{dekegel2021} B.~De~Kegel, N.~Quinn, N.~A.~Thompson, D.~J.~Adams,
C.~J.~Ryan. Comprehensive prediction of robust synthetic lethality between
paralog pairs in cancer cell lines. \emph{Cell Systems} 12(12):1144--1159,
2021.

\bibitem{dekegel2026} B.~De~Kegel, C.~J.~Ryan.
A predicted cancer dependency map for paralog pairs. \emph{bioRxiv}, 2026.

\bibitem{dempster2021} J.~M.~Dempster et al.
Chronos: a cell population dynamics model of CRISPR experiments that
improves inference of gene fitness effects.
\emph{Genome Biology} 22:343, 2021.

\end{thebibliography}

\end{document}
